% \section{Conclusion and Outlook}

% We present a robust framework for longitudinally-informed lesion tracking that secured first place in the MICCAI [\textit{redacted}] challenge
% % autoPET IV Challenge (Task 2)
% and demonstrates state-of-the-art generalization on PanTrack, a completely unseen, out-of-distribution pancreas dataset. We argue that the ``Verified Tracking'' paradigm represents the most clinically viable middle ground between full automation and manual control: it averts catastrophic retrieval failures through minimal human oversight while leveraging deep longitudinal priors to significantly improve delineation accuracy. 

% \noindent Crucially, our approach is designed to integrate seamlessly into standard RECIST 1.1 workflows. By restricting the clinician's role to verifying or quickly correcting a single follow-up point per lesion, this minimal interaction natively resolves complex topological challenges. Vanishing lesions are easily dismissed, and splitting or merging lesions are handled gracefully, as the verification step anchors the network to the clinically relevant sub-component prior to segmentation. 

% \noindent While our current framework focuses on longitudinal tracking and assumes an identified baseline lesion, a prerequisite readily automated by off-the-shelf single-timepoint detectors, it successfully isolates and solves the critical bottleneck of longitudinal correspondence. Furthermore, while we evaluate robustness using rigorously simulated localization errors, our strong zero-shot OOD performance provides a compelling foundation for future prospective clinical reader studies. To foster further research in generalizable tracking, we publicly release the PanTrack dataset, along with our code and model weights.

\section{Conclusion and Outlook}

We present a robust framework for longitudinally-informed lesion tracking that secured first place in the MICCAI autoPET IV challenge and demonstrates state-of-the-art generalization on \textit{PanTrack}, a novel out-of-distribution dataset. We propose ``Verified Tracking'' as a clinically viable middle ground for standard RECIST 1.1 workflows. By requiring clinicians to merely verify or correct a single follow-up point, this paradigm averts catastrophic retrieval failures and gracefully handles complex topological changes like splitting or vanishing lesions. Once anchored, our architecture leverages explicit temporal fusion and synthetic pretraining to fully exploit the baseline appearance, significantly improving delineation accuracy. While our framework currently assumes a known baseline lesion, a step readily automated by off-the-shelf detectors, it successfully isolates and solves the critical bottleneck of longitudinal correspondence. Building on our strong zero-shot OOD performance, future work will focus on prospective clinical reader studies and exploring richer prompt modalities beyond points, such as free-text descriptions of lesion characteristics~\cite{Rokuss_2026_CVPR}. To foster further research in generalizable tracking, we publicly release the \textit{PanTrack} dataset, along with our code and model weights.
